% --- scholar LaTeX fragment --------------------------------------------
% Generated by: scholar sessions export -f table
% Session: noweb-ramsey-v2
% \input this file from the body of a document whose
% preamble loads:
%   \usepackage{longtable}
%   \usepackage{booktabs}
%   \usepackage{array}
% ----------------------------------------------------------------------

\begingroup\footnotesize
\begin{longtable}{@{}%
  >{\raggedright\arraybackslash}p{0.44\textwidth}c%
  >{\raggedright\arraybackslash}p{0.13\textwidth}%
  >{\raggedright\arraybackslash}p{0.26\textwidth}@{}}
\caption{Träffar som rör påståendet, noweb-artikeln (1 citerade, 13 som stöder påståendet och 11 som kvalificerar eller motsäger det, av 93 unika träffar; de 47 angränsande och 20 felträffarna, 1 dubbletter eller andra versioner av citerade källor listas inte). Skälet till varje inkludering står i sista kolumnen; för maskinklassade rader anges modellens konfidens. Databaser: OpenAlex (OA; inte open access), Web of Science (WoS).}\label{tab:sok-noweb}\\
\toprule
Titel & År & Leverantör & Skäl \\
\midrule
\endfirsthead
\multicolumn{4}{@{}l}{\emph{\tablename~\thetable{} (forts.)}}\\
\toprule
Titel & År & Leverantör & Skäl \\
\midrule
\endhead
\midrule \multicolumn{4}{r@{}}{\emph{forts.\ på nästa sida}}\\
\endfoot
\bottomrule
\endlastfoot
Literate programming simplified & 1994 & OA & \emph{citerad}: noweb och dess notation \\
A Perl port of the mathsPIC graphics package & 2000 & OA & stöder påståendet (0.87) \\
A tool for publishing reproducible algorithms \& A reproducible, elegant algorithm for sequential experiments & 2018 & WoS & stöder påståendet (0.85) \\
Comprehension of Literate Programs by Novice and Intermediate Programmers & 2000 & OA & stöder påståendet (0.92) \\
Delaunay graphs by divide and conquer & 1998 & OA & stöder påståendet (0.70) \\
Language-Agnostic Reproducible Data Analysis Using Literate Programming & 2016 & OA & stöder påståendet (0.97) \\
Literate Programming Using Noweb & 1997 & OA & stöder påståendet (0.98) \\
Literate Statistical Practice & 2003 & OA & stöder påståendet (0.98) \\
PROGDOC - A new program documentation system & 2003 & WoS & stöder påståendet (0.98) \\
ProgDOC: A new program documentation system & 2003 & OA & stöder påståendet (0.97) (dubblettpost) \\
Rational points on circles & 1998 & OA & stöder påståendet \\
Survival Analysis in XLISP-Stat. A semiliterate program & 1998 & OA & stöder påståendet (0.76) \\
The CLiP Style of Literate Programming & 2000 & OA & stöder påståendet (0.96) \\
Theory and computation in the study of molecular structure & 2006 & WoS & stöder påståendet (0.82) \\
A WYSIWYG literate programming system (preliminary report) & 1991 & OA & kvalificerar eller motsäger påståendet (0.93) \\
Against the Power of Algorithms Closing, Literate Programming, and Source Code Critique & 2019 & WoS, OA & kvalificerar eller motsäger påståendet (0.70) \\
Application of literate-programming principles for the description of a FORTRAN 90 extension to quaternion arithmetic & 2001 & OA & kvalificerar eller motsäger påståendet (0.94) \\
Electronic Grammars and Reproducible Research & 2012 & OA & kvalificerar eller motsäger påståendet (0.80) \\
FWEB: a literate programming system for Fortran8x & 1990 & OA & kvalificerar eller motsäger påståendet (0.94) \\
Reimagining literate programming & 2009 & OA & kvalificerar eller motsäger påståendet (0.93) \\
Rnoweb: Literate programming with and for R & 2011 & OA & kvalificerar eller motsäger påståendet (0.95) \\
Supporting tailorable program visualisation through literate programming and fisheye views & 2001 & WoS, OA & kvalificerar eller motsäger påståendet (0.86) \\
The literate-programming paradigm & 1991 & OA & kvalificerar eller motsäger påståendet (0.97) \\
Theme-based literate programming & 2002 & WoS & kvalificerar eller motsäger påståendet (0.95) \\
Theme-based literate programming & 2003 & OA & kvalificerar eller motsäger påståendet (0.95) (dubblettpost) \\
\end{longtable}
\endgroup
